\chapter{Conclusions}
\label{chap:conclusions}

Satellite fade nowcasting is useful when it supports timely switching, not
merely accurate intermediate predictions. The conclusions below identify what
the experiments establish about that objective and what remains unresolved.

\section{Contribution and Operational Lessons}

The contribution of this thesis is a decision-level comparison of four fade
nowcasting formulations, rather than a new model architecture. By mapping
different predictive outputs to a common switch reference and making decision
coverage explicit, the study distinguishes accurate intermediate predictions
from useful operational decisions. Its central empirical result is that the
best predictor according to a task's native loss is not necessarily the best
switch controller.

On the external holdout, current-level XGBoost, survival XGBoost-AFT, and raw
PatchTST produced useful but different switching policies: conservative
activation, higher fade coverage, and the strongest autoregressive switch
balance, respectively. These results support the use of signal-only models
for switching, but do not establish a universal winner.
In particular, the retained survival operating point was chosen exploratorily
on test data and still requires validation-only calibration.

The negative results are equally informative. The current-level shapelet
models never activated the switch on the evaluated test grid, while the
long-fade classifiers' strong classification scores did not prevent excessive
switch activity. Neither architectural complexity nor high native accuracy
was sufficient. The practical lesson is to select the predictive target and
assess its decision rule together: an output that is poorly aligned with the
required switch behavior cannot be judged useful from its native score alone.

\section{Answers to the Research Questions}

\paragraph{Can learned models improve switch decisions?}
The experiments support their usefulness: signal-history models can recover
reference switch intervals while limiting unnecessary activation. They do not,
however, establish superiority over a deployed engineering controller.
Well-tuned tabular models remained competitive with neural sequence models,
so greater architectural complexity is not itself evidence of operational
improvement.

\paragraph{Which formulation is most useful?}
The choice depends on the cost of missed fade coverage relative to unnecessary
switch time. Residual-duration survival modeling directly addresses whether an
ongoing fade will persist; current-level regression can support a conservative
policy but estimates a different quantity. Forecasting remains useful when
the future trajectory is also required. Grouped-event classification is better
suited to a high-recall warning than precise switch timing. These differences
make target definition an operational choice, not merely a modeling detail.

\paragraph{Do native metrics explain operational usefulness?}
Not sufficiently. The survival TCN obtained the better integrated Brier score,
whereas XGBoost-AFT produced the better switch sequence at the evaluated
operating point. Similarly, low forecast error did not identify the strongest
autoregressive switch model. Native scores remain necessary to assess the
predicted quantity, but decision timing, active duration, and fragmentation
must be assessed after conversion and post-processing.

\paragraph{Why are a common reference and post-processing necessary?}
They prevent differences in reference construction or decision smoothing from
being mistaken for model improvements. A common timestamp grid is also needed:
task-native scores alone compare decisions over different domains. Reporting
coverage alongside switch quality makes this distinction explicit and shows
where a method would still require a fallback decision in operation.

\section{Limitations}

The main limitation is the size and heterogeneity of the available data. The
external holdout design is stricter than a random split and is appropriate for
testing generalization, but it also means that the final conclusions are based
on one held-out acquisition source. Additional acquisition periods, terminals,
satellite links, and weather conditions would be needed to establish broader
operational robustness.

The thesis deliberately used a signal-only policy. This made the comparison
consistent across files and avoided dependence on columns that were not
available in all datasets. At the same time, it excluded potentially relevant
physical covariates, such as weather measurements, link geometry, or system
state variables. A model that combines signal history with consistently
available physical information could improve anticipation of fades, especially
before the signal crosses the operational threshold.

Another limitation concerns threshold and calibration choices. The switch
reference was built around a \SI{10.0}{dB} threshold and a
\SI{300}{s} operational duration. These values were fixed to make the
experiments comparable, but real deployments may prefer different operating
points. The survival probability threshold used for the retained switch
analysis was exploratory and not a fully validation-calibrated operating
point. Future work should calibrate all decision thresholds on validation data
only and then evaluate the chosen operating point on test data.

Finally, the cross-task comparison required aligning task-native decisions to a
common reference grid. This is necessary for fairness, but it also exposes a
structural issue: different task definitions do not naturally produce
decisions at exactly the same timestamps. The reported coverage values should
therefore be interpreted as part of the result, not as a nuisance. A task that
is accurate only on a small native domain may still be useful, but it cannot
replace a controller that must decide at every operational timestamp.

\section{Future Work}

The first extension should be validation-only calibration of switch-conversion
rules. Autoregressive trajectory thresholds, current-level duration
thresholds, long-fade probability thresholds, and survival probability
thresholds should be tuned on validation data using explicitly chosen
operational costs. This would turn the current comparison from a fixed-rule
benchmark into a calibrated decision study.

Second, the survival formulation should be developed further. XGBoost-AFT
showed that residual-duration modeling is promising, but richer survival
models could improve calibration and event-level timing. Candidate directions
include better discrete-time neural survival models, DeepHit-style competing
risk formulations if multiple event types are introduced, and validation of
survival probability calibration across operational horizons.

Third, the dataset should be expanded and enriched. More external holdout
sources would allow blocked or source-level cross-validation, reducing the risk
that the conclusions depend on one held-out file. If auxiliary variables become
available consistently, future work could compare the current signal-only
framework with a physically enriched framework while preserving the same
no-leakage and switch-evaluation protocol.

Fourth, future experiments should connect the switch metrics to explicit
operational costs. Precision, recall, F1, IoU, active duration, and event
counts describe different aspects of switch behavior, but an operator may
assign asymmetric costs to false activations, missed fades, delayed onsets, and
excessive switch fragmentation. A cost-based decision layer would make the
choice between conservative and high-recall models more transparent.

Fifth, the framework should be moved from offline evaluation toward an online
nowcasting prototype. The current experiments reproduce the information
available at each decision time, but they are still executed as batch analyses.
A deployable system would need streaming ingestion, stateful preprocessing,
incremental gap and quality checks, stateful switch post-processing, latency
measurement, and clear fallback behavior when input quality is poor. Such a
prototype would also make it possible to compare the learned decision rules
against engineering baselines under realistic operating constraints, rather
than only after the acquisition has been processed offline.

The broader conclusion is that satellite fade nowcasting should be treated as a
decision-oriented time-series problem. Accurate forecasts, duration estimates,
and survival curves are valuable only insofar as they support timely and
reliable switch decisions. The framework developed in this thesis provides a
reproducible basis for that evaluation and shows that task formulation,
post-processing, and operational metrics are central components of the
modeling problem, not secondary details.
